\section{Predictive realization and the inherited general mechanism}\label{app:predictive}
This appendix states the precise connection to the mother experiment. It prevents a current task vector from being confused with the minimal state for every executable future test.

Let $\mathcal H$ be the standard Borel disjoint union of finite visible histories, with available actions, stop permissions and all interface coordinates recorded. Choose a countable determining family $(g_j)$ of bounded future tests, containing the required report cylinders and closed under each legal one-step continuation. Relative to a specified prepared law, write
\[
 q(h)=(\E[g_j\mid h])_{j\ge1}
\]
and adjoin the interface coordinates. Countable determination means that equality of these values is equivalent to equality for the entire declared test family. Conditional test values are given jointly Borel versions. Assume explicitly that $q(\mathcal H)$ is Borel and that $q$ has a Borel right inverse on its image. A general Borel surjection is \emph{not} silently assumed to have such a right inverse.

\begin{proposition}[Realization and minimality]\label{prop:quotient}
Under these hypotheses, predictive equivalence $h\sim h'$ if and only if $q(h)=q(h')$ has a standard Borel realization. Its legal report laws and its next predictive state admit jointly Borel versions depending on $(q(h),a,y)$. They agree with the history kernels at almost every report under each actual row. If another statistic $z(h)$ is sufficient for all declared future tests, in the sense that every test value is a Borel function of $z(h)$, then $q$ is a Borel function of $z$. Thus the realization is minimal for those tests and that preparation.
\end{proposition}
\begin{proof}
The image is standard Borel by assumption and its fibers are exactly the equivalence classes. At state $u$ choose the history $\sigma(u)$ given by the Borel right inverse, and define the report and next-state kernels from that history. These are jointly Borel. One-step test closure implies that histories in a fiber have the same integrals of every report cylinder times every next test value. Their report laws are equal, and uniqueness of conditional expectations implies equality of each next-test coordinate almost everywhere under that common law. Countability lets us take a common null set for the coordinates. Hence the next predictive state agrees with the actual history update there. Off that set use the chosen Borel version; no positive-probability continuous singleton is used. Finally the coordinate functions through $z$ assemble into a Borel map into the countable product, with values in the declared image on actual histories, proving minimality.
\end{proof}

The finite raw realizations in Section~\ref{sec:realizations}, with their model fixed and known, give such a representation directly: the reported classical state or pure density matrix determines all future kernel words and is separated by the audit means. No abstract selection assumption is needed there. If the model is unknown and retained under a prior, its posterior may be part of the full predictive state. The current known task response can still be report-recoverable and satisfy \eqref{eq:envelopes}; our proof does not assume that it alone summarizes all future parameter information. Likewise the Gaussian posterior is a history-dependent state, not its current report.

\subsection*{General common-program occupation}
For orientation we record the mechanism retained with complete proofs in~\cite{GTF49}. A finite ordered raw carrier $(E,C,e)$ has positive report maps $L_a(y)$ with $\int eL_a(y)x\,d\lambda_a=e x$. A homogeneous live-state program induces a positive substochastic map $P$ on the product carrier and an exit map $Q$. With prepared vector $a$, its actual live occupation and exit mass are
\begin{equation}\label{eq:occupation-general}
 z_{\rm act}=\sum_{t\ge0}P^ta,\qquad d=Qz_{\rm act}.
\end{equation}
A finite candidate satisfying $z=a+Pz$ need not itself be acquired: iteration gives
\[
 z=\sum_{t=0}^{n-1}P^ta+P^nz.
\]
If the exit mass is one, the decreasing positive limit $h=\lim_nP^nz$ is invariant, $Qh=0$, and $z_{\rm act}=z-h\le z$. Thus nonnegative upper-cost constraints survive removal of inaccessible circulation, but that circulation may not be counted as observed mass. The central regenerative profile is the case in which this actual-action constraint can be solved explicitly, preserving an entire terminal distribution with one live label. The Gaussian proof solves \eqref{eq:occupation-general} for a genuinely informative two-state raw carrier.

For an unknown fixed model, a \emph{single} routing kernel must solve the balance constraints in every model. The retained common-kernel separation theorem sums model and state contributions before maximizing over output labels. Exchanging those operations would permit model-dependent programs. The exact action profile in the present paper has no hidden use of that invalid exchange; its unknown-model extension is instead the explicit common program in Theorem~\ref{thm:central}\textup{(iii)}.

\subsection*{Informative Bayesian compression is not discarded}
The likelihood-weighted identity from~\cite{GTF32} also remains in force. Let affine raw functions $U(y,x)$ and $\ell(y,x)\ge0$ describe a finite future-test family, and let $F=U/\ell$ on positive evidence. For an actual incoming predictive vector $X$ and retained label $S$, put $V=\E[X\mid S]$ and $W=F(Y,V)$. Affine integration before normalization gives $W=\E[X'\mid S,Y]$. For any fresh randomized outgoing label $J$, define $c_{sj}=\E[W\mid S=s,J=j]$ and $v_j=\E[W\mid J=j]$. Two conditional projections yield
\begin{equation}\label{eq:bayesian}
 \E\norm{X'-v_J}^2
 =\E\norm{X'-W}^2
  +\E\norm{W-c_{SJ}}^2+\E\norm{c_{SJ}-v_J}^2.
\end{equation}
Each term is nonnegative. The first is equivalently the Jensen defect of $\int\norm{U(y,x)}^2/\ell(y,x)\,d\lambda(y)$. Zero-evidence reports contribute zero. The exact identity alone does not supply a globally compatible compressed recursion. The prior global-cover and curvature or same-report stability theorems supply sufficient hypotheses for that different task, with their complete proofs retained. Neither task regeneration nor one-step centroid equality is substituted for those hypotheses.


\begin{example}[Local exact cells do not supply lost continuation information]\label{ex:continuation}
Prepare two independent fair signs $U,V$ and reveal both at the first call. At the second call an independent uniform query $Q\in\{1,2\}$ is reported; the protected target is $Z=U$ if $Q=1$ and $Z=V$ otherwise. If the cut retains only $S=U$, then after the query
\[
 W=\E[Z\mid S,Q]=S\one_{\{Q=1\}}.
\]
Three final labels encode $W\in\{-1,0,1\}$ exactly. Every row quantization defect and every pooled-centroid defect in \eqref{eq:bayesian} is zero, but the incoming-information defect is $\E(Z-W)^2=1/2$. More generally, a zero-error answer for both future queries requires at least four cut labels. For randomized encoding, the supports of the cut-label distributions of two different sign pairs must be disjoint: the pairs differ in a coordinate, and that query has positive probability, so a common label would have to decode both signs without access to the discarded pair. Hence four disjoint nonempty supports are necessary. The example concerns the cut, not a free-clock total-state formula. It explains why local centroid equalities cannot replace a global acquired-occupation or continuation certificate.
\end{example}
